% Generated mechanically from the frozen stacks-limits.tex target.
% Do not edit this generated file; edit the bound locale source instead.

% OK, start here.
%

\setcounter{chapter}{101}
\chapter{代数叠的极限}



\phantomsection
\label{stacks-limits-section-phantom}


\section{引言}
\label{stacks-limits-section-introduction}

\noindent
本章汇集与代数叠的极限有关的材料。
关于代数叠和代数空间极限的许多结果由 David Rydh
在 \cite{rydh_approx} 中得到。


\section{约定}
\label{stacks-limits-section-conventions}

\noindent
我们继续采用《叠的性质》第
\ref{stacks-properties-section-conventions} 节中引入的约定和语言滥用。







\section{有限表示态射}
\label{stacks-limits-section-finite-presentation}

\noindent
本节是《空间的极限》第
\ref{spaces-limits-section-finite-presentation} 节的对应版本。
在那里，我们定义了 $\Sch$ 上函子之间的变换保持极限的含义
（建议参看《空间的极限》引理
\ref{spaces-limits-lemma-characterize-relative-limit-preserving}
中的刻画）。在《可表示性判据》第
\ref{criteria-section-limit-preserving} 节中，我们定义了
“在对象上保持极限”的概念。回顾《Artin 公理》第
\ref{artin-section-limits} 节，我们已经定义了 $\Sch$
上的群胚纤维化范畴保持极限的含义。将这些定义结合起来，
得到如下概念。

\begin{definition}
\label{stacks-limits-definition-limit-preserving}
设 $S$ 为概形，并设 $f : \mathcal{X} \to \mathcal{Y}$ 是
$(\Sch/S)_{fppf}$ 上群胚纤维化范畴之间的一个 $1$-态射。
若对 $S$ 上仿射概形的每个有向极限 $U = \lim U_i$，纤维范畴图
$$
\xymatrix{
\colim \mathcal{X}_{U_i} \ar[r] \ar[d]_f & \mathcal{X}_U \ar[d]^f \\
\colim \mathcal{Y}_{U_i} \ar[r] & \mathcal{Y}_U
}
$$
都是 $2$-笛卡儿的，则称 $f$ {\it 保持极限}。
\end{definition}

\begin{lemma}
\label{stacks-limits-lemma-limit-preserving-objects}
设 $S$ 为概形，并设 $f : \mathcal{X} \to \mathcal{Y}$ 是
$(\Sch/S)_{fppf}$ 上群胚纤维化范畴之间的一个 $1$-态射。
若 $f$ 保持极限（定义 \ref{stacks-limits-definition-limit-preserving}），
则 $f$ 在对象上保持极限（《可表示性判据》第
\ref{criteria-section-limit-preserving} 节）。
\end{lemma}

\begin{proof}
% P12-ERR-0001: see ../control/ERRATA_CANDIDATES.jsonl
若对 $U$ 上仿射概形的每个有向极限 $U = \lim U_i$，函子
$$
\colim \mathcal{X}_{U_i} \longrightarrow
(\colim \mathcal{Y}_{U_i}) \times_{\mathcal{Y}_U} \mathcal{X}_U
$$
都是本质满的，则 $f$ 在对象上保持极限。
\end{proof}

\begin{lemma}
\label{stacks-limits-lemma-base-change-limit-preserving}
设 $p : \mathcal{X} \to \mathcal{Y}$ 和
$q : \mathcal{Z} \to \mathcal{Y}$ 是 $(\Sch/S)_{fppf}$
上群胚纤维化范畴之间的 $1$-态射。若
$p : \mathcal{X} \to \mathcal{Y}$ 保持极限，则 $p$ 沿 $q$
的基变换
$p' : \mathcal{X} \times_\mathcal{Y} \mathcal{Z} \to \mathcal{Z}$
也保持极限。
\end{lemma}

\begin{proof}
这是形式推论。设 $U = \lim_{i \in I} U_i$ 是 $S$ 上仿射概形
$U_i$ 的有向极限。对每个 $i$，有
$$
(\mathcal{X} \times_\mathcal{Y} \mathcal{Z})_{U_i} =
\mathcal{X}_{U_i} \times_{\mathcal{Y}_{U_i}} \mathcal{Z}_{U_i}
$$
滤过余极限与范畴的 $2$-纤维积可交换（略去细节），
因此若 $p$ 保持极限，就得到
\begin{align*}
\colim (\mathcal{X} \times_\mathcal{Y} \mathcal{Z})_{U_i}
& =
\colim \mathcal{X}_{U_i} \times_{\colim \mathcal{Y}_{U_i}}
\colim \mathcal{Z}_{U_i} \\
& =
\mathcal{X}_U \times_{\mathcal{Y}_U} \colim \mathcal{Y}_{U_i}
\times_{\colim \mathcal{Y}_{U_i}}
\colim \mathcal{Z}_{U_i} \\
& =
\mathcal{X}_U \times_{\mathcal{Y}_U} \colim \mathcal{Z}_{U_i} \\
& =
\mathcal{X}_U \times_{\mathcal{Y}_U} \mathcal{Z}_U \times_{\mathcal{Z}_U}
\colim \mathcal{Z}_{U_i} \\
& =
(\mathcal{X} \times_\mathcal{Y} \mathcal{Z})_U \times_{\mathcal{Z}_U}
\colim \mathcal{Z}_{U_i}
\end{align*}
这正是所需结论。
\end{proof}

\begin{lemma}
\label{stacks-limits-lemma-composition-limit-preserving}
设 $p : \mathcal{X} \to \mathcal{Y}$ 和
$q : \mathcal{Y} \to \mathcal{Z}$ 是 $(\Sch/S)_{fppf}$
上群胚纤维化范畴之间的 $1$-态射。若 $p$ 和 $q$ 都保持极限，
则复合 $q \circ p$ 也保持极限。
\end{lemma}

\begin{proof}
这是形式推论。设 $U = \lim_{i \in I} U_i$ 是 $S$ 上仿射概形
$U_i$ 的有向极限。若 $p$ 和 $q$ 保持极限，则
\begin{align*}
\colim \mathcal{X}_{U_i}
& =
\mathcal{X}_U \times_{\mathcal{Y}_U} \colim \mathcal{Y}_{U_i} \\
& =
\mathcal{X}_U \times_{\mathcal{Y}_U} \mathcal{Y}_U
\times_{\mathcal{Z}_U} \colim \mathcal{Z}_{U_i} \\
& =
\mathcal{X}_U \times_{\mathcal{Z}_U} \colim \mathcal{Z}_{U_i}
\end{align*}
这正是所需结论。
\end{proof}

\begin{lemma}
\label{stacks-limits-lemma-representable-by-spaces-limit-preserving}
设 $p : \mathcal{X} \to \mathcal{Y}$ 是 $(\Sch/S)_{fppf}$
上群胚纤维化范畴之间的一个 $1$-态射。若 $p$ 可由代数空间表示，
则下列条件等价：
\begin{enumerate}
\item $p$ 保持极限；
\item $p$ 在对象上保持极限；
\item $p$ 局部有限表示（参见《代数叠》定义
\ref{algebraic-definition-relative-representable-property}）。
\end{enumerate}
\end{lemma}

\begin{proof}
由《可表示性判据》引理
\ref{criteria-lemma-representable-by-spaces-limit-preserving}，
我们已经知道 (2) 与 (3) 等价。因此只需证明 (1) 与 (2) 等价。
其中一个方向已见引理 \ref{stacks-limits-lemma-limit-preserving-objects}。
为证另一方向，设 $U = \lim_{i \in I} U_i$ 是 $S$ 上仿射概形
$U_i$ 的有向极限。我们必须证明
$$
\colim \mathcal{X}_{U_i} \longrightarrow
\mathcal{X}_U \times_{\mathcal{Y}_U} \colim \mathcal{Y}_{U_i}
$$
是等价。由假设 (2)，它是本质满的，故只需证明其全忠实性。
由于 $p$ 在纤维范畴上忠实（《代数叠》引理
\ref{algebraic-lemma-criterion-map-representable-spaces-fibred-in-groupoids})
，可见该函子忠实。设 $x_i$ 和 $x'_i$ 是 $\mathcal{X}$ 在
$U_i$ 上的纤维范畴中的对象。上述函子将 $x_i$ 映为
$(x_i|_U, p(x_i), can)$，其中 $can$ 是典范同构
$p(x_i|_U) \to p(x_i)|_U$。于是，设给定态射
$$
(\alpha, \beta_i) : (x_i|_U, p(x_i), can) \longrightarrow
(x'_i|_U, p(x'_i), can)
$$
，它位于本证明第一个陈列箭头右端的范畴中。我们的任务是构造
$i' \geq i$ 以及一个态射
$x_i|_{U_{i'}} \to x'_i|_{U_{i'}}$，使其映到
$(\alpha, \beta_i|_{U_{i'}})$。

\medskip\noindent
令 $y_i = p(x_i)$、$y'_i = p(x'_i)$。由《代数叠》引理
\ref{algebraic-lemma-criterion-map-representable-spaces-fibred-in-groupoids})
，函子
$$ 
X_{y_i} : (\Sch/U_i)^{opp} \to \textit{Sets},\quad
V/U_i \mapsto
% P12-ERR-0002: see ../control/ERRATA_CANDIDATES.jsonl
\{(x, \phi) \mid x \in \Ob(\mathcal{X}_V), \phi : f(x) \to y_i|V\}/\cong
$$
是 $U_i$ 上的代数空间；类似定义的函子 $X_{y'_i}$ 亦然。
由于 (2) 与 (3) 等价，$X_{y'_i}$ 在 $U_i$ 上局部有限表示。
注意，$(x_i, \text{id})$ 和 $(x'_i, \text{id})$
分别定义了 $X_{y_i}$ 和 $X_{y'_i}$ 的 $U_i$-值点。
存在函子变换
$$
\beta_i : X_{y_i} \to X_{y'_i},\quad
(x/V, \phi) \mapsto (x/V, \beta_i|_V \circ \phi)
$$
；换言之，这是 $U_i$ 上代数空间之间的一个态射。我们断言
$$
\xymatrix{
U \ar[d] \ar[rr] & & U_i \ar[d]^{(x'_i, \text{id})} \\
U_i \ar[r]^{(x_i, \text{id})} & X_{y_i} \ar[r]^{\beta_i} & X_{y'_i}
}
$$
交换。事实上，这等价于函子 $X_{y'_i}$ 定义中的两对
$(x_i|_U, \beta_i|_U)$ 与 $(x'_i|_U, \text{id})$ 同构；
而态射 $\alpha : x_i|_U \to x'_i|_U$ 恰好给出这样的同构。
反向推理可见，若能找到 $i' \geq i$，使得图
$$
\xymatrix{
U_{i'} \ar[d] \ar[rr] & & U_i \ar[d]^{(x'_i, \text{id})} \\
U_i \ar[r]^{(x_i, \text{id})} & X_{y_i} \ar[r]^{\beta_i} & X_{y'_i}
}
$$
交换，那么就得到同构 $x_i|_{U_{i'}} \to x'_i|_{U_{i'}}$，
它解决了上一段提出的问题。另一方面，对角态射
$$
\Delta : X_{y'_i} \to X_{y'_i} \times_{U_i} X_{y'_i}
$$
局部有限表示（《空间的态射》引理
\ref{spaces-morphisms-lemma-diagonal-morphism-finite-type})
。因此，$U \to U_i$ 使到 $X_{y'_i}$ 的两个态射相等这一事实，
意味着对某个 $i' \geq i$，态射 $U_{i'} \to U_i$
已使这两个态射相等；参见《空间的极限》命题
\ref{spaces-limits-proposition-characterize-locally-finite-presentation}。
\end{proof}

\begin{lemma}
\label{stacks-limits-lemma-limit-preserving-diagonal}
设 $p : \mathcal{X} \to \mathcal{Y}$ 是 $(\Sch/S)_{fppf}$
上群胚纤维化范畴之间的一个 $1$-态射。下列条件等价：
\begin{enumerate}
\item 对角态射
$\Delta : \mathcal{X} \to \mathcal{X} \times_\mathcal{Y} \mathcal{X}$
保持极限；
\item 对 $S$ 上仿射概形的每个有向极限 $U = \lim U_i$，函子
$$
\colim \mathcal{X}_{U_i} \longrightarrow
\mathcal{X}_U \times_{\mathcal{Y}_U} \colim \mathcal{Y}_{U_i}
$$
是全忠实的。
\end{enumerate}
特别地，若 $p$ 保持极限，则 $\Delta$ 也保持极限。
\end{lemma}

\begin{proof}
设 $U = \lim U_i$ 是 $S$ 上仿射概形的有向极限。我们断言函子
$$
\colim \mathcal{X}_{U_i} \longrightarrow
\mathcal{X}_U \times_{\mathcal{Y}_U} \colim \mathcal{Y}_{U_i}
$$
全忠实，当且仅当函子
$$
\colim \mathcal{X}_{U_i} \longrightarrow
\mathcal{X}_U \times_{(\mathcal{X} \times_\mathcal{Y} \mathcal{X})_U}
\colim (\mathcal{X} \times_\mathcal{Y} \mathcal{X})_{U_i}
$$
是等价。这将证明本引理。由于
$(\mathcal{X} \times_\mathcal{Y} \mathcal{X})_U =
\mathcal{X}_U \times_{\mathcal{Y}_U} \mathcal{X}_U$
并且
$(\mathcal{X} \times_\mathcal{Y} \mathcal{X})_{U_i} =
\mathcal{X}_{U_i} \times_{\mathcal{Y}_{U_i}} \mathcal{X}_{U_i}$
，这纯粹是一个范畴论断言；下一段将予讨论。

\medskip\noindent
设 $\mathcal{I}$ 为滤过指标范畴，并设
$(\mathcal{C}_i)$ 和 $(\mathcal{D}_i)$ 是 $\mathcal{I}$
上的群胚系统。设 $p : (\mathcal{C}_i) \to (\mathcal{D}_i)$
为 $\mathcal{I}$ 上群胚系统之间的映射。再设有群胚函子
$p : \mathcal{C} \to \mathcal{D}$ 以及函子
$f : \colim \mathcal{C}_i \to \mathcal{C}$ 和
$g : \colim \mathcal{D}_i \to \mathcal{D}$，它们构成交换图
$$
\xymatrix{
\colim \mathcal{C}_i \ar[d]_p \ar[r]_f & \mathcal{C} \ar[d]^p \\
\colim \mathcal{D}_i \ar[r]^g & \mathcal{D}
}
$$
我们断言，
$$
A : \colim \mathcal{C}_i \longrightarrow
\mathcal{C} \times_\mathcal{D} \colim \mathcal{D}_i
$$
全忠实，当且仅当函子
$$
B : \colim \mathcal{C}_i \longrightarrow
\mathcal{C}
\times_{\Delta, \mathcal{C} \times_\mathcal{D} \mathcal{C}, f \times_g f}
\colim (\mathcal{C}_i \times_{\mathcal{D}_i} \mathcal{C}_i)
$$
是等价。令 $\mathcal{C}' = \colim \mathcal{C}_i$、
$\mathcal{D}' = \colim \mathcal{D}_i$。由于 $2$-纤维积与滤过余极限
可交换，$A$ 和 $B$ 分别成为函子
$$
A' : \mathcal{C}' \to \mathcal{C} \times_\mathcal{D} \mathcal{D}'
\quad\text{且}\quad
B' : \mathcal{C}' \longrightarrow
\mathcal{C}
\times_{\Delta, \mathcal{C} \times_\mathcal{D} \mathcal{C}, f \times_g f}
(\mathcal{C}' \times_{\mathcal{D}'} \mathcal{C}')
$$
因此只需证明：若
$$
\xymatrix{
\mathcal{C}' \ar[d]_p \ar[r]_f & \mathcal{C} \ar[d]^p \\
\mathcal{D}' \ar[r]^g & \mathcal{D}
}
$$
是群胚的交换图，则 $A'$ 全忠实，当且仅当 $B'$ 是等价。
这可由《范畴》引理
\ref{categories-lemma-fully-faithful-diagonal-equivalence}
（取平凡的，即单点的基范畴）推出，因为
$$
\mathcal{C}
\times_{\Delta, \mathcal{C} \times_\mathcal{D} \mathcal{C}, f \times_g f}
(\mathcal{C}' \times_{\mathcal{D}'} \mathcal{C}') =
\mathcal{C}'
\times_{A', \mathcal{C} \times_\mathcal{D} \mathcal{D}', A'}
\mathcal{C}'
$$
证明完毕。
\end{proof}

\begin{lemma}
\label{stacks-limits-lemma-locally-finite-presentation-limit-preserving}
设 $S$ 为概形，$\mathcal{X}$ 为 $S$ 上的代数叠。
若 $\mathcal{X} \to S$ 局部有限表示，则 $\mathcal{X}$
在《Artin 公理》定义 \ref{artin-definition-limit-preserving}
的意义下保持极限（等价地，态射 $\mathcal{X} \to S$ 保持极限）。
\end{lemma}

\begin{proof}
选取某个概形 $U$ 及一个光滑满射 $U \to \mathcal{X}$。
则 $U \to S$ 局部有限表示，参见《叠的态射》第
\ref{stacks-morphisms-section-finite-presentation} 节。
由《代数叠》引理 \ref{algebraic-lemma-stack-presentation}，
对代数空间中的某个光滑群胚 $(U, R, s, t, c)$，可写成
$\mathcal{X} = [U/R]$。由于 $U$ 在 $S$ 上局部有限表示，
代数空间 $R$ 在 $S$ 上也局部有限表示。
回顾，$[U/R]$ 是将如下群胚纤维化范畴叠化而得到的
$(\Sch/S)_{fppf}$ 上的群胚叠：它在 $T$ 上的纤维范畴是群胚
$(U(T), R(T), s, t, c)$。由于作为函子，$U$ 和 $R$ 都保持极限
（《空间的极限》命题
\ref{spaces-limits-proposition-characterize-locally-finite-presentation})
，该群胚纤维化范畴保持极限。因此，只需证明 fppf 叠化保留
保持极限这一性质。事实确实如此（提示：使用《拓扑》引理
\ref{topologies-lemma-limit-fppf-topology}).
不过，利用本情形中叠化的具体描述，我们在下面给出直接证明。

\medskip\noindent
设 $T = \lim T_\lambda$ 是 $S$ 上仿射概形的有向极限。
我们必须证明函子
$$
\colim [U/R]_{T_\lambda} \longrightarrow [U/R]_T
$$
是范畴等价。先证明该函子本质满。设 $x \in \Ob([U/R]_T)$。
《空间中的群胚》引理
\ref{spaces-groupoids-lemma-quotient-stack-objects}
给出了范畴 $[U/R]_T$ 的描述。特别地，$x$ 对应于一个 fppf 覆盖
$\{T_i \to T\}_{i \in I}$ 以及相对于该覆盖的一个
$[U/R]$-下降数据 $(u_i, r_{ij})$。细化此覆盖后，可设它是仿射概形
$T$ 的标准 fppf 覆盖。由《拓扑》引理
\ref{topologies-lemma-limit-fppf-topology}
，可选取某个 $\lambda$ 和标准 fppf 覆盖
$\{T_{\lambda, i} \to T_\lambda\}_{i \in I}$，使其到 $T$
的基变换等于 $\{T_i \to T\}_{i \in I}$。
对每个 $i$，增大 $\lambda$ 后，可找到
$u_{\lambda, i} : T_{\lambda, i} \to U$，使其与
$T_i \to T_{\lambda, i}$ 的复合为给定态射 $u_i$
（此处使用 $U$ 保持极限）。类似地，对每个 $i, j$，增大
$\lambda$ 后，可找到
$r_{\lambda, ij} : T_{\lambda, i} \times_{T_\lambda} T_{\lambda, j} \to R$，
使其与 $T_{ij} \to T_{\lambda, ij}$ 的复合为给定态射 $r_{ij}$
（此处使用 $R$ 保持极限）。进一步增大 $\lambda$ 后，可设
$$
s \circ r_{\lambda, ij} = u_{\lambda, i} \circ \text{pr}_0
\quad\text{且}\quad
t \circ r_{\lambda, ij} = u_{\lambda, j} \circ \text{pr}_1,
$$
并且
$$
c \circ (r_{\lambda, jk} \circ \text{pr}_{12}, r_{\lambda, ij}
\circ \text{pr}_{01}) = r_{\lambda, ik} \circ \text{pr}_{02}.
$$
换言之，可设 $(u_{\lambda, i}, r_{\lambda, ij})$ 是相对于覆盖
$\{T_{\lambda, i} \to T_\lambda\}_{i \in I}$ 的一个
$[U/R]$-下降数据。于是得到 $T_\lambda$ 上 $[U/R]$ 的相应对象，
其到 $T$ 的拉回与 $x$ 同构，正合所需。
全忠实性的证明完全相同，只需使用《空间中的群胚》引理
\ref{spaces-groupoids-lemma-quotient-stack-objects}
% P12-ERR-0003: see ../control/ERRATA_CANDIDATES.jsonl
给出的 $[U/T]$ 的纤维范畴中态射的描述。
\end{proof}

\begin{proposition}
\label{stacks-limits-proposition-characterize-locally-finite-presentation}
\begin{reference}
这是 \cite[Lemma 2.3.15]{Emerton-Gee} 的一个特例。
\end{reference}
设 $f : \mathcal{X} \to \mathcal{Y}$ 为代数叠之间的态射。
下列条件等价：
\begin{enumerate}
\item $f$ 保持极限；
\item $f$ 在对象上保持极限；
\item $f$ 局部有限表示。
\end{enumerate}
\end{proposition}

\begin{proof}
假设 (3) 成立。设 $T = \lim T_i$ 是仿射概形的有向极限。
考虑函子
$$
\colim \mathcal{X}_{T_i} \longrightarrow
\mathcal{X}_T \times_{\mathcal{Y}_T} \colim \mathcal{Y}_{T_i}
$$
设 $(x, y_i, \beta)$ 是右端的对象，即
$x \in \Ob(\mathcal{X}_T)$、$y_i \in \Ob(\mathcal{Y}_{T_i})$，
且 $\beta : f(x) \to y_i|_T$ 是 $\mathcal{Y}_T$ 中的态射。
于是可将 $(x, y_i, \beta)$ 视为 $T$ 上代数叠
$\mathcal{X}_{y_i} = \mathcal{X} \times_{\mathcal{Y}, y_i} T_i$
的一个对象。由于 $\mathcal{X}_{y_i} \to T_i$ 是 $f$ 的基变换，
故它局部有限表示；由引理
\ref{stacks-limits-lemma-locally-finite-presentation-limit-preserving}，
它保持极限。这意味着，对某个 $i' \geq i$，
$(x, y_i, \beta)$ 来自 $T_{i'}$ 上的对象。展开定义即知，
$(x, y_i, \beta)$ 属于所陈列函子的本质像。
换言之，该函子本质满；也就是说，$f$ 在对象上保持极限。
现在将此结论应用于 $f$ 的对角态射 $\Delta$。由《叠的态射》引理
\ref{stacks-morphisms-lemma-diagonal-morphism-finite-type}
，态射 $\Delta$ 局部有限表示。因此上述论证表明 $\Delta$
在对象上保持极限。由引理
\ref{stacks-limits-lemma-representable-by-spaces-limit-preserving}，
$\Delta$ 保持极限。再由引理 \ref{stacks-limits-lemma-limit-preserving-diagonal}，
上述陈列函子全忠实。它又已被证明本质满，故是等价，从而 (1) 成立。

\medskip\noindent
蕴含关系 (1) $\Rightarrow$ (2) 显然。现假设 (2)。
选取概形 $V$ 及光滑满射 $V \to \mathcal{Y}$。
由《可表示性判据》引理
\ref{criteria-lemma-base-change-limit-preserving}
，基变换 $\mathcal{X} \times_\mathcal{Y} V \to V$
在对象上保持极限。选取概形 $U$ 及光滑满射
$U \to \mathcal{X} \times_\mathcal{Y} V$。
由于光滑态射局部有限表示，证明的第一部分说明
$U \to \mathcal{X} \times_\mathcal{Y} V$ 保持极限。
由《可表示性判据》引理
\ref{criteria-lemma-composition-limit-preserving}
，复合 $U \to V$ 在对象上保持极限。因此 $U \to V$
局部有限表示，参见《可表示性判据》引理
\ref{criteria-lemma-representable-by-spaces-limit-preserving}。
这恰好是 $f$ 局部有限表示的条件，参见《叠的态射》定义
\ref{stacks-morphisms-definition-locally-finite-presentation}。
\end{proof}




\section{性质的下降}
\label{stacks-limits-section-descent}

\noindent
本节是《极限》第 \ref{limits-section-descent} 节的对应版本。

\begin{situation}
\label{stacks-limits-situation-descent}
设 $Y = \lim_{i \in I} Y_i$ 是一个代数空间有向系统的极限，
其转移态射均为仿射态射。假设对每个 $i \in I$，$X_i$
% P12-ERR-0004: see ../control/ERRATA_CANDIDATES.jsonl
都是拟紧且拟分离的。另选定一个元素 $0 \in I$。
\end{situation}

\begin{lemma}
\label{stacks-limits-lemma-eventually-separated}
在情形 \ref{stacks-limits-situation-descent} 中，设
$\mathcal{X}_0 \to Y_0$ 是从代数叠到 $Y_0$ 的态射，
并设 $\mathcal{X}_0$ 拟紧且拟分离。若
$Y \times_{Y_0} \mathcal{X}_0 \to Y$ 是分离态射，则对所有充分大的
$i \in I$，$Y_i \times_{Y_0} \mathcal{X}_0 \to Y_i$ 都是分离态射。
\end{lemma}

\begin{proof}
记 $\mathcal{X} = Y \times_{Y_0} \mathcal{X}_0$、
$\mathcal{X}_i = Y_i \times_{Y_0} \mathcal{X}_0$。
选取仿射概形 $U_0$ 及光滑满射 $U_0 \to \mathcal{X}_0$。
令 $U = Y \times_{Y_0} U_0$、$U_i = Y_i \times_{Y_0} U_0$。
则 $U$ 和 $U_i$ 均仿射，且 $U \to \mathcal{X}$ 和
$U_i \to \mathcal{X}_i$ 均为光滑满射。令
$R_0 = U_0 \times_{\mathcal{X}_0} U_0$、
$R = Y \times_{Y_0} R_0$、$R_i = Y_i \times_{Y_0} R_0$。
于是 $R = U \times_\mathcal{X} U$ 且
$R_i = U_i \times_{\mathcal{X}_i} U_i$。

\medskip\noindent
在此记号下，若 $\mathcal{X} \to Y$ 分离，则
$R \to U \times_Y U$ 是对角态射
$\mathcal{X} \to \mathcal{X} \times_Y \mathcal{X}$
沿 $U \times_Y U \to \mathcal{X} \times_Y \mathcal{X}$ 的基变换，
故为固有态射。反之，若 $R_i \to U_i \times_{Y_i} U_i$ 固有，
则 $\mathcal{X}_i \to Y_i$ 分离，因为
$U_i \times_{Y_i} U_i \to \mathcal{X}_i \times_{Y_i} \mathcal{X}_i$
是光滑满射；参见《叠的性质》引理
\ref{stacks-properties-lemma-check-property-covering}。
注意，$R_0 \to U_0 \times_{Y_0} U_0$ 局部有限型，且 $R_0$
拟紧、拟分离。由《空间的极限》引理
\ref{spaces-limits-lemma-eventually-proper}，对充分大的 $i$，
$R_i \to U_i \times_{Y_i} U_i$ 固有，证明完毕。
\end{proof}







\section{相对对象的下降}
\label{stacks-limits-section-descending-relative}

\noindent
本节是《空间的极限》第
\ref{spaces-limits-section-descending-relative} 节的对应版本。

\begin{lemma}
\label{stacks-limits-lemma-descend-a-stack-down}
设 $I$ 为有向集，$(X_i, f_{ii'})$ 为以 $I$ 为指标的代数空间逆系统。
假设
\begin{enumerate}
\item 态射 $f_{ii'} : X_i \to X_{i'}$ 均为仿射态射；
\item 空间 $X_i$ 均拟紧且拟分离。
\end{enumerate}
令 $X = \lim X_i$。若 $\mathcal{X}$ 是 $X$ 上有限表示的代数叠，
则存在 $i \in I$ 及 $X_i$ 上有限表示的代数叠 $\mathcal{X}_i$，
使得作为 $X$ 上的代数叠有
$\mathcal{X} \cong \mathcal{X}_i \times_{X_i} X$。
\end{lemma}

\begin{proof}
由《叠的态射》定义
\ref{stacks-morphisms-definition-locally-finite-presentation}
，态射 $\mathcal{X} \to X$ 拟紧、局部有限表示且拟分离。
由于 $X$ 拟紧且 $\mathcal{X} \to X$ 拟紧，$\mathcal{X}$ 也拟紧
（《叠的态射》定义
\ref{stacks-morphisms-definition-quasi-compact}）。
因此，可找到仿射概形 $U$ 及光滑满射
$U \to \mathcal{X}$（《叠的性质》引理
\ref{stacks-properties-lemma-quasi-compact-stack}）。
令 $R = U \times_\mathcal{X} U$。由《代数叠》引理
\ref{algebraic-lemma-stack-presentation}，得到 $X$ 上代数空间中的
光滑群胚 $(U, R, s, t, c)$，使得 $\mathcal{X} = [U/R]$。
由于 $\mathcal{X} \to X$ 拟分离且 $X$ 拟分离，
$\mathcal{X}$ 拟分离（《叠的态射》引理
\ref{stacks-morphisms-lemma-composition-separated}）。
故 $R \to U \times U$ 拟紧且拟分离（《叠的态射》引理
\ref{stacks-morphisms-lemma-fibre-product-after-map}），
从而 $R$ 是拟分离、拟紧的代数空间。另一方面，$U \to X$
局部有限表示，因此 $R \to X$ 也局部有限表示
（因为 $s : R \to U$ 光滑，故局部有限表示）。
于是 $(U, R, s, t, c)$ 是 $X$ 上有限表示代数空间范畴中的群胚对象。
由《空间的极限》引理
\ref{spaces-limits-lemma-descend-finite-presentation}
，存在某个 $i$ 及 $X_i$ 上代数空间中的群胚
$(U_i, R_i, s_i, t_i, c_i)$，其到 $X$ 的拉回与
$(U, R, s, t, c)$ 同构。增大 $i$ 后，可设 $s_i$ 和 $t_i$
均光滑，参见《空间的极限》引理
\ref{spaces-limits-lemma-descend-smooth}。商叠
$\mathcal{X}_i = [U_i/R_i]$ 是代数叠（《代数叠》定理
\ref{algebraic-theorem-smooth-groupoid-gives-algebraic-stack}）。

\medskip\noindent
由《空间中的群胚》引理
\ref{spaces-groupoids-lemma-quotient-stack-functorial}，
有态射 $[U/R] \to [U_i/R_i]$。我们断言，把它与态射
$[U/R] \to X$ 和 $[U_i/R_i] \to X_i$
（《空间中的群胚》引理
\ref{spaces-groupoids-lemma-quotient-stack-arrows}）结合起来，
得到同构（即等价）
$$
[U/R] \longrightarrow [U_i/R_i] \times_{X_i} X
$$
相应的映射
$$
[U/_{\!p}R] \longrightarrow [U_i/_{\!p}R_i] \times_{X_i} X
$$
在《空间中的群胚》公式
(\ref{spaces-groupoids-equation-quotient-stack}) 所述的“群胚预层”
层次上是同构。因此，该断言由叠化与纤维积可交换这一事实推出，
参见《叠》引理
\ref{stacks-lemma-stackification-fibre-product-fibred-categories}。
\end{proof}




\section{有限型对象嵌入有限表示对象}
\label{stacks-limits-section-finite-type-closed-in-finite-presentation}

\noindent
本节是《空间的极限》第
\ref{spaces-limits-section-finite-type-closed-in-finite-presentation}
节的对应版本。

\begin{lemma}
\label{stacks-limits-lemma-finite-type-closed-in-finite-presentation}
设 $f : \mathcal{X} \to Y$ 是从代数叠到代数空间的态射。假设：
\begin{enumerate}
\item $f$ 有限型且拟分离；
\item $Y$ 拟紧且拟分离。
\end{enumerate}
则存在有限表示态射 $f' : \mathcal{X}' \to Y$，以及 $Y$
上代数叠之间的闭浸入 $\mathcal{X} \to \mathcal{X}'$。
\end{lemma}

\begin{proof}
将 $Y = \lim_{i \in I} Y_i$ 写成以有向集 $I$ 为指标的代数空间极限，
其中转移态射为仿射态射，且 $Y_i$ 为诺特代数空间；参见
《空间的极限》命题 \ref{spaces-limits-proposition-approximate}。
我们将使用《空间的极限》第
\ref{spaces-limits-section-finite-type-quasi-separated} 节的材料。

\medskip\noindent
选取一个表示 $\mathcal{X} = [U/R]$，并以
$(U, R, s, t, c, e, i)$ 表示相应的 $Y$ 上代数空间群胚。
可以并且确实设 $U$ 仿射。于是 $U$、$R$ 和
$R \times_{s, U, t} R$ 都是 $Y$ 上有限型、拟分离的代数空间。
% P12-ERR-0005: see ../control/ERRATA_CANDIDATES.jsonl
有两个态射 $s, t : R \to U$，三个态射
$c : R \times_{s, U, t} R \to R$,
$\text{pr}_1 : R \times_{s, U, t} R \to R$,
$\text{pr}_2 : R \times_{s, U, t} R \to R$,
，一个态射 $e : U \to R$，以及最后一个态射 $i : R \to R$。
这些态射满足一组公理；详见《群胚》第
\ref{groupoids-section-groupoids} 节。

\medskip\noindent
根据《空间的极限》注
\ref{spaces-limits-remark-finite-type-gives-well-defined-system}
，可找到 $i_0 \in I$ 以及逆系统
\begin{enumerate}
\item $(U_i)_{i \geq i_0}$,
\item $(R_i)_{i \geq i_0}$,
\item $(T_i)_{i \geq i_0}$
\end{enumerate}
，它们位于 $(Y_i)_{i \geq i_0}$ 上，并满足
$U = \lim_{i \geq i_0} U_i$,
$R = \lim_{i \geq i_0} R_i$, and
$R \times_{s, U, t} R = \lim_{i \geq i_0} T_i$
，而且存在系统态射
\begin{enumerate}
\item $(s_i)_{i \geq i_0} : (R_i)_{i \geq i_0} \to (U_i)_{i \geq i_0}$,
\item $(t_i)_{i \geq i_0} : (R_i)_{i \geq i_0} \to (U_i)_{i \geq i_0}$,
\item $(c_i)_{i \geq i_0} : (T_i)_{i \geq i_0} \to (R_i)_{i \geq i_0}$,
\item $(p_i)_{i \geq i_0} : (T_i)_{i \geq i_0} \to (R_i)_{i \geq i_0}$,
\item $(q_i)_{i \geq i_0} : (T_i)_{i \geq i_0} \to (R_i)_{i \geq i_0}$,
\item $(e_i)_{i \geq i_0} : (U_i)_{i \geq i_0} \to (R_i)_{i \geq i_0}$,
\item $(i_i)_{i \geq i_0} : (R_i)_{i \geq i_0} \to (R_i)_{i \geq i_0}$
\end{enumerate}
使得
$s = \lim_{i \geq i_0} s_i$,
$t = \lim_{i \geq i_0} t_i$,
$c = \lim_{i \geq i_0} c_i$,
$\text{pr}_1 = \lim_{i \geq i_0} p_i$,
$\text{pr}_2 = \lim_{i \geq i_0} q_i$,
$e = \lim_{i \geq i_0} e_i$, and
$i = \lim_{i \geq i_0} i_i$.
由《空间的极限》引理
\ref{spaces-limits-lemma-morphism-good-diagram-smooth}
，可设 $s_i$ 和 $t_i$ 均光滑（这可能需要增大 $i_0$）。
由《空间的极限》引理
\ref{spaces-limits-lemma-morphism-good-diagram-flat}
，可设对所有 $i \geq i_0$，由 $s$ 和 $R \to R_i$ 给出的映射
$R \to U \times_{U_i, s_i} R_i$，以及由 $t$ 和 $R \to R_i$
给出的映射 $R \to U \times_{U_i, t_i} R_i$，都是同构。
由《空间的极限》引理
\ref{spaces-limits-lemma-good-diagram-fibre-product}，可设各图
$$
\xymatrix{
T_i \ar[r]_{q_i} \ar[d]_{p_i} & R_i \ar[d]^{t_i} \\
R_i \ar[r]^{s_i} & U_i
}
$$
均为笛卡儿图。《空间的极限》引理
\ref{spaces-limits-lemma-morphism-good-diagram} 的唯一性于是保证：
当 $i$ 充分大时，$s_i, t_i, c_i, e_i, i_i$
满足上述态射 $s, t, c, e, i$ 之间的关系。固定这样的一个 $i$。

\medskip\noindent
于是 $(U_i, R_i, s_i, t_i, c_i, e_i, i_i)$ 是 $Y_i$
上代数空间中的光滑群胚。因此 $\mathcal{X}_i = [U_i/R_i]$
是代数叠（《代数叠》定理
\ref{algebraic-theorem-smooth-groupoid-gives-algebraic-stack}).
群胚态射
$$
(U, R, s, t, c, e, i) \to (U_i, R_i, s_i, t_i, c_i, e_i, i_i)
$$
位于 $Y \to Y_i$ 上，并确定交换图
$$
\xymatrix{
\mathcal{X} \ar[d] \ar[r] & \mathcal{X}_i \ar[d] \\
Y \ar[r] & Y_i
}
$$
（《空间中的群胚》引理
\ref{spaces-groupoids-lemma-quotient-stack-functorial}).
我们断言态射 $\mathcal{X} \to Y \times_{Y_i} \mathcal{X}_i$
是闭浸入。该断言将完成证明，因为按构造，代数叠
$\mathcal{X}_i \to Y_i$ 有限表示。为证此断言，注意左图
$$
\xymatrix{
U \ar[d] \ar[r] & U_i \ar[d] \\
\mathcal{X} \ar[r] & \mathcal{X}_i
}
\quad\quad
\xymatrix{
U \ar[d] \ar[r] & Y \times_{Y_i} U_i \ar[d] \\
\mathcal{X} \ar[r] & Y \times_{Y_i} \mathcal{X}_i
}
$$
由《空间中的群胚》引理
\ref{spaces-groupoids-lemma-criterion-fibre-product}
及上述结果可知是笛卡儿图，故右侧交换图也是笛卡儿图。
根据《空间的极限》引理
\ref{spaces-limits-lemma-limit-from-good-diagram} 中逆系统
$(U_i)$ 的构造，$U \to Y \times_{Y_i} U_i$ 是闭浸入；
而 $Y \times_{Y_i} U_i \to Y \times_{Y_i} \mathcal{X}_i$
是光滑满射。结合《叠的性质》引理
\ref{stacks-properties-lemma-check-immersion-covering}，
即得所需结论。
\end{proof}

\noindent
对于分离代数叠，也有如下版本。

\begin{lemma}
\label{stacks-limits-lemma-separated-closed-in-finite-presentation}
设 $f : \mathcal{X} \to Y$ 是从代数叠到代数空间的态射。假设：
\begin{enumerate}
\item $f$ 有限型且分离；
\item $Y$ 拟紧且拟分离。
\end{enumerate}
则存在分离的有限表示态射 $f' : \mathcal{X}' \to Y$，
以及 $Y$ 上代数叠之间的闭浸入
$\mathcal{X} \to \mathcal{X}'$。
\end{lemma}

\begin{proof}
首先完全按照引理
\ref{stacks-limits-lemma-finite-type-closed-in-finite-presentation}
证明中的步骤（并沿用其中的记号），把嵌入
$\mathcal{X} \to \mathcal{X}'$ 构造成态射
$\mathcal{X} \to \mathcal{X}' = Y \times_{Y_i} \mathcal{X}_i$，
其中 $\mathcal{X}_i = [U_i/R_i]$。因此只需证明当 $i$
充分大时，$\mathcal{X}_i \to Y_i$ 分离。换言之，只需证明当 $i$
充分大时，$\mathcal{X}_i \to \mathcal{X}_i \times_{Y_i} \mathcal{X}_i$
固有。由于态射
$U_i \times_{Y_i} U_i \to \mathcal{X}_i \times_{Y_i} \mathcal{X}_i$
是光滑满射，且
$R_i = \mathcal{X}_i
\times_{\mathcal{X}_i \times_{Y_i} \mathcal{X}_i} U_i \times_{Y_i} U_i$
，只需证明态射
$(s_i, t_i) : R_i \to U_i \times_{Y_i} U_i$
在 $i$ 充分大时固有；参见《叠的性质》引理
\ref{stacks-properties-lemma-check-property-covering}。
下一段证明这一点。

\medskip\noindent
注意，$U \times_Y U \to Y$ 拟分离且有限型。因此可使用
《空间的极限》注
\ref{spaces-limits-remark-finite-type-gives-well-defined-system}
的构造，找到 $i_1 \in I$ 以及逆系统 $(V_i)_{i \geq i_1}$，
使 $U \times_Y U = \lim_{i \geq i_1} V_i$。
由《空间的极限》引理
\ref{spaces-limits-lemma-good-diagram-fibre-product}，
当 $i$ 充分大时，将该构造的函子性用于投影
$U \times_Y U \to U$，得到闭浸入
$$
V_i \to U_i \times_{Y_i} U_i
$$
（这里略有不符，因为严格说来应以 $Y \to Y_i$ 的概形论像替换
$Y_i$，但显然这不改变该纤维积。）另一方面，由《空间的极限》引理
\ref{spaces-limits-lemma-morphism-good-diagram-proper}
，把函子性用于固有态射
$(s, t) : R \to U \times_Y U$（此处使用 $\mathcal{X}$ 分离），
便得到当 $i$ 充分大时为固有态射的 $R_i \to V_i$。
复合这些态射可得：对所有充分大的 $i$，态射
$R_i \to U_i \times_{Y_i} U_i$ 固有。《空间的极限》注
\ref{spaces-limits-remark-finite-type-gives-well-defined-system}
% P12-ERR-0005: see ../control/ERRATA_CANDIDATES.jsonl
中构造的函子性表明，当 $i$ 充分大时，此态射就是
$(s_i, t_i)$，证明完毕。
\end{proof}






\section{泛闭态射}
\label{stacks-limits-section-universally-closed}

\noindent
本节是《空间的极限》第
\ref{spaces-limits-section-universally-closed} 节的对应版本。

\begin{lemma}
\label{stacks-limits-lemma-separate}
设 $g : Z \to Y$ 为仿射概形之间的态射，并设
$f : \mathcal{X} \to Y$ 为代数叠之间的拟紧态射。
设 $z \in Z$，且 $T \subset |\mathcal{X} \times_Y Z|$
是满足 $z \not \in \Im(T \to |Z|)$ 的闭子集。
若 $\mathcal{X}$ 拟紧，则存在开邻域 $V \subset Z$（含 $z$）、交换图
$$
\xymatrix{
V \ar[d] \ar[r]_a & Z' \ar[d]^b \\
Z \ar[r]^g & Y,
}
$$ 
% P12-ERR-0006: see ../control/ERRATA_CANDIDATES.jsonl
以及闭子集 $T' \subset |X \times_Y Z'|$，使得
\begin{enumerate}
\item $Z'$ 是 $Y$ 上有限表示的仿射概形；
\item 令 $z' = a(z)$，则 $z' \not \in \Im(T' \to |Z'|)$；
\item $T$ 在 $|\mathcal{X} \times_Y V|$ 中的逆像经由
$|\mathcal{X} \times_Y V| \to |\mathcal{X} \times_Y Z'|$
映入 $T'$。
\end{enumerate}
\end{lemma}

\begin{proof}
我们从概形态射的相应结果推出本结论。由于 $\mathcal{X}$ 拟紧，
可选取仿射概形 $W$ 及光滑满射 $W \to \mathcal{X}$。
令 $T_W \subset |W \times_Y Z|$ 为 $T$ 的逆像。
则 $z$ 不属于 $T_W$ 的像。由概形情形（《极限》引理
\ref{limits-lemma-separate}），可找到含 $z$ 的开邻域
$V \subset Z$、概形的交换图
$$
\xymatrix{
V \ar[d] \ar[r]_a & Z' \ar[d]^b \\
Z \ar[r]^g & Y,
}
$$
以及闭子集 $T' \subset |W \times_Y Z'|$，使得
\begin{enumerate}
\item $Z'$ 是 $Y$ 上有限表示的仿射概形；
\item 令 $z' = a(z)$，则 $z' \not \in \Im(T' \to |Z'|)$；
\item $T_1 = T_W \cap |W \times_Y V|$ 经由
$|W \times_Y V| \to |W \times_Y Z'|$ 映入 $T'$。
\end{enumerate}
交换图
$$
\xymatrix{
W \times_Y Z \ar[d] &
W \times_Y V \ar[l] \ar[rr]_{a_1} \ar[d]_c & &
W \times_Y Z' \ar[d]^q \\
\mathcal{X} \times_Y Z &
\mathcal{X} \times_Y V \ar[l] \ar[rr]^{a_2} & &
\mathcal{X} \times_Y Z'
}
$$
的两个方块均为笛卡儿方块，且竖直映射为光滑满射，因而尤其是开映射。
考察左侧方块可见，
$T_1 = T_W \cap |W \times_Y V|$ 是
$T_2 = T \cap |\mathcal{X} \times_Y V|$ 在 $c$ 下的逆像。
由《叠的性质》引理
\ref{stacks-properties-lemma-points-cartesian}，得到
$a_1(T_1) = q^{-1}(a_2(T_2))$。
由《拓扑》引理 \ref{topology-lemma-open-morphism-quotient-topology}，
得到
$$
q^{-1}\left(\overline{a_2(T_2)}\right) =
\overline{q^{-1}(a_2(T_2))} =
\overline{a_1(T_1)} \subset T'
$$
由于 $q$ 满射，$\overline{a_2(T_2)} \to |Z'|$ 的像不含 $z'$，
因为对 $T'$ 亦然。
因此，取上述含 $Z', V, a, b$ 的图以及闭子集
$\overline{a_2(T_2)} \subset |\mathcal{X} \times_Y Z'|$，
即得到本引理所述问题的解。
\end{proof}

\begin{lemma}
\label{stacks-limits-lemma-test-universally-closed}
设 $f : \mathcal{X} \to \mathcal{Y}$ 为代数叠之间的拟紧态射。
下列条件等价：
\begin{enumerate}
\item $f$ 泛闭；
% P12-ERR-0007: see ../control/ERRATA_CANDIDATES.jsonl
\item 对每个局部有限表示态射 $Z \to \mathcal{Y}$，其中 $Z$
为仿射概形，映射 $|\mathcal{X} \times_Y Z| \to |Z|$ 是闭映射；
\item 存在概形 $V$ 及光滑满射 $V \to \mathcal{Y}$，使得
$|\mathbf{A}^n \times (\mathcal{X} \times_\mathcal{Y} V)|
\to |\mathbf{A}^n \times V|$ 对所有 $n \geq 0$ 都是闭映射。
\end{enumerate}
\end{lemma}

\begin{proof}
(1) 蕴含 (2) 显然。

\medskip\noindent
假设 (2)。选取一个由仿射概形不交并而成的概形 $V$，以及光滑满射
$V \to \mathcal{Y}$。为证明 $f$ 泛闭，只需证明 $f$ 的基变换
$\mathcal{X} \times_\mathcal{Y} V \to V$ 泛闭；参见《叠的态射》引理
\ref{stacks-morphisms-lemma-universally-closed-local}。
注意，对该基变换，条件 (2) 仍成立。因此在证明 (2) 蕴含 (1) 时，
可设 $Y = \mathcal{Y}$ 是仿射概形。

\medskip\noindent
假设 (2)，并设 $\mathcal{Y} = Y$ 为仿射概形。若 $f$ 不泛闭，
则存在 $Y$ 上仿射概形 $Z$，使
$|\mathcal{X} \times_Y Z| \to |Z|$ 不是闭映射；参见《叠的态射》引理
\ref{stacks-morphisms-lemma-universally-closed-local}。
这意味着存在闭子集 $T \subset |\mathcal{X} \times_Y Z|$，
使 $\Im(T \to |Z|)$ 不闭。选取 $z \in |Z|$，使其属于 $T$
之像的闭包但不属于该像。应用引理 \ref{stacks-limits-lemma-separate}，
得到开邻域 $V \subset Z$、交换图
$$
\xymatrix{
V \ar[d] \ar[r]_a & Z' \ar[d]^b \\
Z \ar[r]^g & Y,
}
$$
以及闭子集 $T' \subset |\mathcal{X} \times_Y Z'|$，使得
\begin{enumerate}
\item $Z'$ 是 $Y$ 上有限表示的仿射概形；
\item 令 $z' = a(z)$，则 $z' \not \in \Im(T' \to |Z'|)$；
\item $T$ 在 $|\mathcal{X} \times_Y V|$ 中的逆像经由
$|\mathcal{X} \times_Y V| \to |\mathcal{X} \times_Y Z'|$
映入 $T'$。
\end{enumerate}
我们断言 $z'$ 属于 $\Im(T' \to |Z'|)$ 的闭包。这意味着
$|\mathcal{X} \times_Y Z'| \to |Z'|$ 不是闭映射，
与假设 (2) 矛盾。换言之，该断言证明了 (2) 蕴含 (1)。
为证明此断言，考察交换图
$$
\xymatrix{
\mathcal{X} \times_Y Z \ar[d] &
\mathcal{X} \times_Y V \ar[l] \ar[d] \ar[r] &
\mathcal{X} \times_Y Z' \ar[d] \\
Z & V \ar[l] \ar[r]^a & Z'
}
$$
令 $T_V \subset |\mathcal{X} \times_Y V|$ 为 $T$ 的逆像。
由《叠的性质》引理 \ref{stacks-properties-lemma-points-cartesian}，
$T_V$ 在 $|V|$ 中的像等于 $T$ 在 $|Z|$ 中之像的逆像。
由于 $z$ 属于 $T \to |Z|$ 之像的闭包，且 $|V| \to |Z|$ 是开映射，
$z$ 属于 $T_V \to |V|$ 之像的闭包。又因 $T_V$ 在
$|\mathcal{X} \times_Y Z'|$ 中的像包含于 $|T'|$，
立即得到 $z' = a(z)$ 属于 $T'$ 之像的闭包。

\medskip\noindent
(1) 蕴含 (3) 显然。设 $V \to \mathcal{Y}$ 如 (3) 所述。
% P12-ERR-0008: see ../control/ERRATA_CANDIDATES.jsonl
若能证明 $\mathcal{X} \times_Y V \to V$ 泛闭，则由《叠的态射》引理
\ref{stacks-morphisms-lemma-universally-closed-local}，$f$ 泛闭。
因此，只需证明：若 $f : \mathcal{X} \to \mathcal{Y}$
是代数叠之间的拟紧态射，$\mathcal{Y} = Y$ 是概形，且对所有 $n$，
$|\mathbf{A}^n \times \mathcal{X}| \to |\mathbf{A}^n \times Y|$
均为闭映射，则 $f$ 满足 (2)。设 $Z \to Y$ 局部有限表示，
其中 $Z$ 为仿射概形。我们必须证明映射
$|\mathcal{X} \times_Y Z| \to |Z|$ 是闭映射。
由于 $Y$ 是概形、$Z$ 仿射且 $Z \to Y$ 局部有限表示，
可找到浸入 $Z \to \mathbf{A}^n \times Y$；参见《态射》引理
\ref{morphisms-lemma-quasi-affine-finite-type-over-S}。考虑笛卡儿图
$$
\vcenter{
\xymatrix{
\mathcal{X} \times_Y Z \ar[d] \ar[r] &
\mathbf{A}^n \times \mathcal{X} \ar[d] \\
Z \ar[r] & \mathbf{A}^n \times Y
}
}
\quad
\begin{matrix}
\text{诱导出} \\
\text{Cartesian 方块}
\end{matrix}
\quad
\vcenter{
\xymatrix{
|\mathcal{X} \times_Y Z| \ar[d] \ar[r] &
|\mathbf{A}^n \times \mathcal{X}| \ar[d] \\
|Z| \ar[r] & |\mathbf{A}^n \times Y|
}
}
$$
，其中水平箭头都是到局部闭子集的同胚（《叠的性质》引理
\ref{stacks-properties-lemma-subspace-induced-topology}）。
% P12-ERR-0009: see ../control/ERRATA_CANDIDATES.jsonl
因此，$|X \times_Y Z|$ 的每个闭子集 $T$ 都是
$|\mathbf{A}^n \times Y|$ 中某个闭子集 $T'$ 的拉回。
假设说明 $T'$ 在 $|\mathbf{A}^n \times X|$ 中的像闭，
故 $T$ 在 $|Z|$ 中的像闭，正合所需。
\end{proof}
